% فهرست منابع
{\singlespacing
\begin{latin}
\begin{thebibliography}{99}
\phantomsection
\addcontentsline{toc}{chapter}{فهرست منابع}

\bibitem{raissi2019}
M. Raissi, P. Perdikaris, and G. E. Karniadakis, “Physics-informed neural
networks: A deep learning framework for solving forward and inverse problems
involving nonlinear partial differential equations,” Journal of Computational
Physics, vol. 378, pp. 686–707, 2019.

\bibitem{lagaris1998}
I. E. Lagaris, A. Likas, and D. I. Fotiadis, “Artificial neural networks for
solving ordinary and partial differential equations,” IEEE Transactions on
Neural Networks, vol. 9, no. 5, pp. 987–1000, 1998.

\bibitem{raissi2017a}
M. Raissi, P. Perdikaris, and G. E. Karniadakis, “Machine learning of linear
differential equations using Gaussian processes,” Journal of Computational
Physics, vol. 348, pp. 683–693, 2017.

\bibitem{raissi2018}
M. Raissi and G. E. Karniadakis, “Hidden physics models: Machine learning of
nonlinear partial differential equations,” Journal of Computational Physics,
vol. 357, pp. 125–141, 2018.

\bibitem{karpatne2017}
A. Karpatne, G. Atluri, J. Faghmous, M. Steinbach, A. Banerjee, A. Ganguly,
S. Shekhar, N. Samatova, and V. Kumar, “Theory-guided data science: A new
paradigm for scientific discovery from data,” IEEE Transactions on Knowledge
and Data Engineering, vol. 29, no. 10, pp. 2318–2331, 2017.

\bibitem{schmidt2009}
M. Schmidt and H. Lipson, “Distilling free-form natural laws from experimental
data,” Science, vol. 324, pp. 81–85, 2009.

\bibitem{sindy2016}
S. L. Brunton, J. L. Proctor, and J. N. Kutz, “Discovering governing equations
from data by sparse identification of nonlinear dynamical systems,” Proceedings
of the National Academy of Sciences, vol. 113, no. 15, pp. 3932–3937, 2016.

\bibitem{cuomo2022}
S. Cuomo, V. S. Di Cola, F. Giampaolo, G. Rozza, M. Raissi, and F. Piccialli,
“Scientific machine learning through physics-informed neural networks: Where
we are and what’s next,” Journal of Scientific Computing, vol. 92, article 88,
2022.

\bibitem{cybenko1989}
G. Cybenko, “Approximation by superpositions of a sigmoidal function,”
Mathematics of Control, Signals and Systems, vol. 2, no. 4, pp. 303–314, 1989.

\bibitem{rumelhart1986}
D. E. Rumelhart, G. E. Hinton, and R. J. Williams, “Learning representations
by back-propagating errors,” Nature, vol. 323, pp. 533–536, 1986.

\bibitem{goodfellow2016}
I. Goodfellow, Y. Bengio, and A. Courville, Deep Learning. MIT Press, 2016.

\bibitem{kingma2015}
D. P. Kingma and J. Ba, “Adam: A method for stochastic optimization,” in
Proceedings of the International Conference on Learning Representations
(ICLR), 2015.

\bibitem{liu1989}
D. C. Liu and J. Nocedal, “On the limited memory BFGS method for large scale
optimization,” Mathematical Programming, vol. 45, pp. 503–528, 1989.

\bibitem{glorot2010}
X. Glorot and Y. Bengio, “Understanding the difficulty of training deep
feedforward neural networks,” in Proceedings of AISTATS, pp. 249–256, 2010.

\bibitem{baydin2018}
A. G. Baydin, B. A. Pearlmutter, A. A. Radul, and J. M. Siskind, “Automatic
differentiation in machine learning: A survey,” Journal of Machine Learning
Research, vol. 18, pp. 1–43, 2018.

\bibitem{leveque2007}
R. J. LeVeque, Finite Difference Methods for Ordinary and Partial Differential
Equations. SIAM, 2007.

\bibitem{evans2010}
L. C. Evans, Partial Differential Equations, 2nd ed. American Mathematical
Society, 2010.

\bibitem{hopf1950}
E. Hopf, “The partial differential equation $u_t+uu_x=\mu u_{xx}$,”
Communications on Pure and Applied Mathematics, vol. 3, pp. 201–230, 1950.

\bibitem{benton1972}
E. R. Benton and G. W. Platzman, “A table of solutions of the one-dimensional
Burgers equation,” Quarterly of Applied Mathematics, vol. 30, pp. 195–212,
1972.

\bibitem{trefethen2000}
L. N. Trefethen, Spectral Methods in MATLAB. SIAM, 2000.

\bibitem{kassam2005}
A.-K. Kassam and L. N. Trefethen, “Fourth-order time-stepping for stiff PDEs,”
SIAM Journal on Scientific Computing, vol. 26, no. 4, pp. 1214–1233, 2005.

\bibitem{butcher2016}
J. C. Butcher, Numerical Methods for Ordinary Differential Equations, 3rd ed.
Wiley, 2016.

\bibitem{abadi2016}
M. Abadi et al., “TensorFlow: Large-scale machine learning on heterogeneous
systems,” in Proceedings of OSDI, 2016.

\bibitem{jax2018}
J. Bradbury et al., “JAX: Composable transformations of Python+NumPy
programs,” 2018. [Online]. Available: github.com/google/jax

\end{thebibliography}
\end{latin}
}
